%!TEX program = xelatex
\documentclass[11pt,a4paper]{article}
\usepackage[margin=2.2cm]{geometry}
\usepackage{amsmath,amssymb,amsthm,mathtools}
\usepackage{bm}
\usepackage{enumitem}
\usepackage{booktabs}
\usepackage{xcolor}
\usepackage{hyperref}
\usepackage{fancyhdr}
\usepackage{xeCJK}

\hypersetup{
    colorlinks=true,
    linkcolor=blue!60!black,
    citecolor=blue!60!black,
    urlcolor=blue!60!black
}

\pagestyle{fancy}
\fancyhf{}
\rhead{\small Qiuzhen QE Probability \& Statistics --- Spring 2025}
\lhead{\small Official Qualifying Exam}
\rfoot{\small Page \thepage}

\DeclareMathOperator*{\argmax}{arg\,max}
\DeclareMathOperator*{\argmin}{arg\,min}
\DeclareMathOperator{\Var}{Var}
\DeclareMathOperator{\E}{\mathbb{E}}
\DeclareMathOperator{\Pbb}{\mathbb{P}}
\DeclareMathOperator{\Cov}{Cov}

\setlist[itemize]{topsep=3pt,itemsep=2pt,parsep=0pt}
\setlist[enumerate]{topsep=3pt,itemsep=2pt,parsep=0pt}

\title{\textbf{QUALIFYING EXAM: PROBABILITY \& STATISTICS}\\[0.5em]\Large SPRING 2025}
\date{}
\author{}

\begin{document}

\maketitle
\vspace{-3em}

\section*{Instructions}
\begin{itemize}
    \item There are 11 problems in this exam. You need to choose 8 of them to solve. If you select more than 8, only the first 8 that you have worked on will be graded. Note that 4 of the problems are worth 15 points each and the rest 10 points each.
    \item You must follow all the rules of exam taking. Misconducts will be subject to proper disciplinary actions by the Center.
    \item You must provide all necessary details for full credits. A final answer with no or little explanation/derivation, even if correct, receives a minimal credit.
\end{itemize}

\noindent\rule{\linewidth}{0.4pt}

\begin{enumerate}[label=\textbf{\arabic*.}]
    \item \textbf{(10 points)} Let $\{U_n\}$ be i.i.d.\ uniform on $[0, 1]$. For $n \ge 2$, define
    \[
    X_n \coloneqq \min\{U_1, \dots, U_n\}, \quad Y_n \coloneqq \max\{U_1, \dots, U_n\}.
    \]
    \begin{enumerate}[label=(\arabic*)]
        \item Derive the joint distribution of $(X_n, Y_n)$.
        \item Calculate $\E [X_n \mid \sigma(Y_n)]$.
    \end{enumerate}

    \item \textbf{(10 points)} Suppose $(X, Y, Z)$ is a 3D mean-zero Gaussian vector with covariance matrix
    \[
    \begin{pmatrix} 1 & \rho & \tau \\ \rho & 1 & 0 \\ \tau & 0 & 1 \end{pmatrix}.
    \]
    \begin{enumerate}[label=(\arabic*)]
        \item Show that $\rho^2 + \tau^2 \le 1$.
        \item Show that there exist real numbers $a, b, c$ and a standard Gaussian random variable $W$ independent of $(Y, Z)$ such that $X = aY + bZ + cW$, and derive $a, b, c$.
    \end{enumerate}

    \item \textbf{(15 points)} Suppose $\{\xi_j\}$ are i.i.d.\ Bernoulli random variables with $\Pbb[\xi_1 = 1] = \Pbb[\xi_1 = -1] = 1/2$. Set $S_n = \sum_{j=1}^n \xi_j$ and define $\tau = \min\{n : S_n = -a \text{ or } b\}$ for $a, b \in \mathbb{Z}_{>0}$.
    \begin{enumerate}[label=(\arabic*)]
        \item Calculate $\E[e^{\lambda S_n}]$ for $\lambda \in \mathbb{R}$.
        \item Calculate $\E[s^\tau]$ for $s \in (0, 1)$.
    \end{enumerate}

    \item \textbf{(10 points)} Suppose $\{\xi_j\}$ are i.i.d.\ Bernoulli random variables with $\Pbb[\xi_1 = 1] = \Pbb[\xi_1 = -1] = 1/2$. Set $X_n = \sum_{j=1}^n 2^{j-1} \xi_j$.
    \begin{enumerate}[label=(\arabic*)]
        \item Define $T = \min\{j \ge 1 : \xi_j = +1\}$. Derive the law of $T$ and calculate $\E[T]$.
        \item Show that $\{X_n\}$ is a martingale and calculate $\E[X_T]$.
    \end{enumerate}

    \item \textbf{(10 points)} Let $\{B_t\}_{t \ge 0}$ be a 1D Brownian motion with $B_0 = 0$. For $a \in (0, \infty)$, define $T_a \coloneqq \inf\{t \ge 0 : |B_t| = a\}$.
    \begin{enumerate}[label=(\arabic*)]
        \item Prove that $T_a < \infty$ a.s.
        \item Prove that $T_a$ and $\mathbf{1}_{\{B_{T_a} = a\}}$ are independent.
    \end{enumerate}

    \item \textbf{(15 points)} Suppose $X$ and $Y$ are i.i.d.\ with $\E[X] = 0$ and $\Var(X) = 1$. Suppose further that
    \[
    \E [(X - Y)^2 \mid \sigma(X + Y)] = 2.
    \]
    Show that $X$ and $Y$ are independent standard Gaussian random variables.

    \item \textbf{(10 points)} Let $X_1, \dots, X_n$ ($n \ge 2$) be i.i.d.\ Uniform $U(0, \theta)$ where $0 < \theta < \infty$.
    \begin{enumerate}[label=(\alph*)]
        \item Under squared error loss $l(\theta, \delta(\mathbf{X})) = (\delta(\mathbf{X}) - \theta)^2$, derive the generalized Bayes rule for the improper prior $\pi(\theta) = 1, 0 < \theta < \infty$.
        \item Is the Bayes rule derived in (a) consistent for $\theta$?
    \end{enumerate}

    \item \textbf{(10 points)} Let $X_1, \dots, X_n$ be i.i.d.\ Gamma$(2025, \theta)$ with $\theta > 0$ unknown.
    \begin{enumerate}[label=(\alph*)]
        \item Consider the testing problem $H_0 : \theta = 1$ vs. $H_1 : \theta \neq 1$. Derive the UMP level $\alpha = 0.05$ test if it exists. If it does not exist, argue why.
        \item Derive the MLE for $\frac{1}{\theta}$, and find its asymptotic distribution.
    \end{enumerate}

    \item \textbf{(15 points)} Let $(X_1, Y_1), \dots, (X_n, Y_n)$ be i.i.d.\ bivariate random vectors with correlation $\rho$. Let $\hat{\rho}$ be the sample correlation coefficient.
    \begin{enumerate}[label=(\alph*)]
        \item Assuming $\E|X_i|^4 < \infty$ and $\E|Y_i|^4 < \infty$, prove that $\sqrt{n}(\hat{\rho} - \rho) \xrightarrow{(d)} \mathcal{N}(0, c)$ and identify $c$ in terms of moments.
        \item Further assume $(X_i, Y_i)$ follows a bivariate normal distribution. Find a minimal sufficient statistic for $\rho$. Is it complete? Explain.
    \end{enumerate}

    \item \textbf{(10 points)} Let $X_1, \dots, X_n$ be i.i.d.\ $\mathcal{N}(0, \sigma^2)$ with $\sigma > 0$ unknown. Let $W = \frac{1}{n} \sum_{i=1}^n |X_i|$.
    \begin{enumerate}[label=(\alph*)]
        \item Find an unbiased estimator of $\sigma$ based on $W$ and identify its limiting distribution.
        \item Is it admissible under squared error loss? Justify your answer.
    \end{enumerate}

    \item \textbf{(15 points)} Let $X_1, \dots, X_n$ be i.i.d.\ Exponential with density $f(x; \lambda, \alpha) = \lambda e^{-\lambda(x - \alpha)} I_{\{x \ge \alpha\}}$, where $\lambda > 0$ and $\alpha \in \mathbb{R}$.
    \begin{enumerate}[label=(\alph*)]
        \item Find the MLE of $(\lambda, \alpha)$, denoted by $(\hat{\lambda}, \hat{\alpha})$.
        \item What distribution does $\frac{\sqrt{n}(\hat{\lambda} - \lambda)}{n(\hat{\alpha} - \alpha)}$ converge to?
    \end{enumerate}
\end{enumerate}

\end{document}
